\section{Introduction}
Computational notebooks such as Jupyter have become central tools in scientific computing, teaching, and data analysis. Their interactive structure—mixing narrative text, code, and visual output—makes them highly suitable for inquiry-based learning in programming education~\cite{perkel2018jupyter}. Yet despite their pedagogical advantages, instructors often lack visibility into how students actually work within these environments. Students, on the other hand, often have little opportunity to reflect on their own workflow, error patterns, or problem-solving behavior. Research in programming education consistently shows that learners repeat similar types of mistakes, struggle to identify the sources of their errors, and often adopt inefficient trial-and-error strategies. At the same time, teachers typically only see final code submissions, which provide limited insight into the learners' process~\cite{robinson2022error}. This creates a gap between the correctness of the final code and the path the student took to reach it. Understanding this process is crucial to improve teaching interventions and help students develop their metacognitive skills.

Learning dashboards have been shown to enhance students' self-awareness, foster reflection, and offer instructors actionable insights into learner behavior. However, the majority of existing dashboards focus on high-level learning management system data rather than fine-grained programming behavior within computational notebooks~\cite{paulsen2024learning}. To bridge this gap, it is first necessary to develop tools that translate raw notebook interactions into meaningful visualizations and, crucially, to evaluate whether students perceive such tools as useful and easy to integrate into their workflow.

This project addresses this gap by designing a JupyterLab dashboard extension capable of tracking detailed notebook interactions, such as cell-level execution histories and error patterns. The primary objective of this study is to evaluate the usability, perceived usefulness, and cognitive workload of this newly developed dashboard from a learner's perspective. By focusing on the user experience (UX) and technology acceptance, this research aims to determine if the proposed system provides a viable foundation for supporting self-reflection in programming education.
To evaluate the system, an exploratory study was conducted with 22 students. Participants engaged with the dashboard while solving Python exercises and subsequently provided feedback through standardized instruments, including the System Usability Scale (SUS), the Technology Acceptance Model (TAM), and the NASA Task Load Index (NASA-TLX)~\cite{germanupa_sus, davis1989technology, hart1988development}. This approach allows for a comprehensive assessment of whether the dashboard meets the needs of students and identifies potential barriers to its adoption. The findings contribute to the field of learning analytics by providing empirical evidence on the design and acceptance of fine-grained, student-centered reflection tools in computational environments.
